\section{曲面的同调}
\label{ch12-sec75}

尽管我们已经成功地得到了许多曲面的基本群的表示, 但目前仍不拟追究这些结论到底对我们有什么用. 从先前的计算中我们得出什么结论呢? 譬如, 双重环面与三重环面是否同胚? 答案并不明显. 我们知道, 并没有一个有效的方法从两个群的表示来判明这两个群是否同构. 倘若考虑阿贝尔群 $\pi_{1}/[\pi_{1},\ \pi_{1}]$, 其中 $\pi_{1}=\pi_{1}(X,\ x_{0})$, 事情就好办多了. 因为这时有一些熟知的不变量可供应用. 本节我们就此进行探讨.

我们知道, 如果 $X$ 是一个道路连通空间, 并且 $\alpha$ 是 $X$ 中从 $x_{0}$ 到 $x_{1}$ 的一条道路, 则存在一个从以 $x_{0}$ 为基点的基本群到以 $x_{1}$ 为基点的基本群的同构 $\widehat{\alpha}$, 但这一同构依赖于道路 $\alpha$ 的选取. 然而对于群 $\pi_{1}/[\pi_{1},\ \pi_{1}]$ 却有一个更强一点的结论. 在这种情况下, 由 $\alpha$ 诱导的从以 $x_{0}$ 为基点的"基本群的阿贝尔化"到以 $x_{1}$ 为基点的相应基本群的同构便与道路 $\alpha$ 的选取无关.

为了验证以上事实, 我们只需证明: 如果 $\alpha$ 和 $\beta$ 是从 $x_{0}$ 到 $x_{1}$ 的两条道路, 则道路 $g=\alpha*\bar{\beta}$ 诱导出从 $\pi_{1}/[\pi_{1},\ \pi_{1}]$ 到自身的恒等同构. 证明这一点很容易. 如果 $[f]\in\pi_{1}(X,\ x_{0})$, 则有
\[
\widehat{g}[f]=[\bar{g}*f*g]=[g]^{-1}*[f]*[g].
\]
若考虑阿贝尔群 $\pi_{1}/[\pi_{1},\ \pi_{1}]$ 中的陪集, 便可见 $\widehat{g}$ 诱导出恒等映射.

\begin{definition}[第一个同调群]
\label{definition:ch12-first-homology}
若 $X$ 是一个道路连通空间, 令
\[
H_{1}(X)=\pi_{1}(X,\ x_{0})/[\pi_{1}(X,\ x_{0}),\ \pi_{1}(X,\ x_{0})].
\]
$H_{1}(X)$ 称为是 $X$ 的\textbf{第一个同调群}(first homology group). 我们在此表达式中省略了基点, 这是由于在两个不同基点的基本群的阿贝尔化之间存在唯一的一个道路诱导同构.
\end{definition}

在以后学习代数拓扑时, 你将会发现那里对 $H_{1}(X)$ 有完全不同的定义. 事实上, 那时对于所有 $n\geqslant 0$, $H_{n}(X)$ 都有定义, 并且称之为 $X$ 的\textbf{同调群}(homology group). 这些阿贝尔群都是 $X$ 的拓扑不变量, 在运用代数理论中的结论去处理拓扑问题时起着重要作用. W. Hurewicz 的一个定理给出了这些群与 $X$ 的同伦群之间的联系. 这个定理告诉我们道路连通空间 $X$ 的第一个同调群 $H_{1}(X)$ 同构于 $X$ 的基本群的阿贝尔化. 这个定理使我们能够将空间的基本群的阿贝尔化定义为第一个同调群.

为了对前面讨论的曲面计算 $H_{1}(X)$, 我们需要以下结论.

\begin{theorem}[商群与阿贝尔化的交换性]
\label{theorem:ch12-abelianization-exchange}
设 $F$ 是一个群, $N$ 是 $F$ 的一个正规子群, $q\colon F\to F/N$ 是投射. 投射同态 $p\colon F\to F/[F,\ F]$ 诱导一个同构
\[
\phi\colon q(F)/[q(F),q(F)]\longrightarrow p(F)/p(N).
\]
\end{theorem}

粗略地讲, 这个定理告诉我们: 用 $N$ 分解 $F$ 再将其商群阿贝尔化, 与先对 $F$ 阿贝尔化再用 $N$ 在阿贝尔化过程中的像来分解它, 效果相同.

\par
\begin{figure}[htbp]
\centering
\includegraphics[width=0.85\textwidth]{image/fig-75-diagram.png}
\end{figure}

\begin{proof}
我们有下图所示的各种投射同态 $p$, $q$, $r$, $s$, 其中 $q(F)=F/N$, $p(F)=F/[F,\ F]$.


由于 $r\circ p$ 将 $N$ 映为 $1$, 它诱导出一个同态 $u\colon q(F)\to p(F)/p(N)$. 因为 $p(F)/p(N)$ 是阿贝尔群, 同态 $u$ 诱导 $q(F)/[q(F),\ q(F)]$ 的一个同态 $\phi$. 另一方面, 由于 $s\circ q$ 将 $F$ 映入一个阿贝尔群, 它诱导出一个同态 $v\colon p(F)\to q(F)/[q(F),\ q(F)]$. 由于 $s\circ q$ 将 $N$ 映为 $1$, 因此 $v\circ p$ 也将 $N$ 映为 $1$, 从而 $v$ 诱导 $p(F)/p(N)$ 的一个同态 $\psi$.

同态 $\phi$ 可以如下描述: 给定群 $q(F)/[q(F),\ q(F)]$ 中的一个元素 $y$, 选取 $F$ 的一个元素 $x$ 使得 $s(q(x))=y$, 则 $\phi(y)=r(p(x))$. 对同态 $\psi$ 可以给出相仿的描述. 由此得到 $\phi$ 和 $\psi$ 是互逆的.
\end{proof}

\begin{corollary}[自由群商群的阿贝尔化]
\label{corollary:ch12-free-group-abelianization}
设 $F$ 是以 $\alpha_{1},\ \cdots,\ \alpha_{n}$ 为自由生成元的自由群, $N$ 为 $F$ 中含有元素 $x\in F$ 的最小正规子群, $G=F/N$, $p\colon F\to F/[F,\ F]$ 为投射. 则 $G/[G,\ G]$ 同构于商群 $F/[F,\ F]$, 它是以 $p(\alpha_{1}),\ \cdots,\ p(\alpha_{n})$ 为基的自由阿贝尔群关于由 $p(x)$ 所生成的子群的商群.
\end{corollary}

\begin{proof}
由于 $N$ 是由 $x$ 及其所有共轭元所生成的, 所以群 $p(N)$ 是由 $p(x)$ 所生成的. 因此由 \cref{theorem:ch12-abelianization-exchange} 可得到本推论.
\end{proof}

\begin{theorem}[环面 $n$-重连通和的同调]
\label{theorem:ch12-torus-sum-homology}
若 $X$ 是环面的 $n$-重连通和, 则 $H_{1}(X)$ 是秩为 $2n$ 的自由阿贝尔群.
\end{theorem}

\begin{proof}
根据 \cref{corollary:ch12-free-group-abelianization}, \cref{theorem:ch12-n-fold-torus-pi1} 蕴涵着 $H_{1}(X)$ 同构于集合 $\alpha_{1},\ \beta_{1},\ \cdots,\ \alpha_{n},\ \beta_{n}$ 上的自由阿贝尔群 $F'$ 关于由元素 $[\alpha_{1},\ \beta_{1}]\cdots[\alpha_{n},\ \beta_{n}]$ 生成的子群的商群, 其中 $[\alpha,\ \beta]=\alpha\beta\alpha^{-1}\beta^{-1}$. 由于群 $F'$ 是阿贝尔群, 这个元素等于单位元.
\end{proof}

\begin{theorem}[射影平面 $m$-重连通和的同调]
\label{theorem:ch12-projective-sum-homology}
若 $X$ 是射影平面的 $m$-重连通和, 则 $H_{1}(X)$ 的挠子群 $T(X)$ 的阶为 $2$, 并且 $H_{1}(X)/T(X)$ 是秩为 $m-1$ 的自由阿贝尔群.
\end{theorem}

\begin{proof}
根据 \cref{corollary:ch12-free-group-abelianization}, \cref{theorem:ch12-m-fold-projective-pi1} 蕴涵着 $H_{1}(X)$ 同构于集合 $\alpha_{1},\ \cdots,\ \alpha_{m}$ 上的自由阿贝尔群 $F'$ 关于由 $(\alpha_{1})^{2},\ \cdots,\ (\alpha_{m})^{2}$ 生成的子群的商群. 如果用加法表示群运算 (这做法在处理阿贝尔群时常用), 这个子群由元素 $2(\alpha_{1}+\cdots+\alpha_{m})$ 生成. 改变群 $F'$ 的基. 若设 $\beta=\alpha_{1}+\cdots+\alpha_{m}$, 则元素 $\alpha_{1},\ \cdots,\ \alpha_{m-1},\ \beta$ 构成 $F'$ 的一个基, 并且 $F'$ 的每一个元素可以由它们唯一地表示. 群 $H_{1}(X)$ 同构于 $\alpha_{1},\ \cdots,\ \alpha_{m-1},\ \beta$ 上的自由阿贝尔群关于由 $2\beta$ 生成的子群的商群. 也就是说, $H_{1}(X)$ 同构于 $m$-重笛卡儿积 $\mathbb{Z}\times\cdots\times\mathbb{Z}$ 关于子群 $0\times\cdots\times 0\times 2\,\mathbb{Z}$ 的商群. 定理得证.
\end{proof}

\begin{theorem}[紧致曲面的非同胚列表]
\label{theorem:ch12-nonhomeomorphic-list}
设 $T_{n}$ 和 $P_{m}$ 分别表示环面的 $n$-重连通和与射影平面的 $m$-重连通和. 则曲面 $S^{2},\ T_{1},\ T_{2},\ \cdots,\ P_{1},\ P_{2},\ \cdots$ 中的任何两个都不同胚.
\end{theorem}

\begin{exercise}
\label{exercise:ch12-sec75-01}
计算 $H_{1}(P^{2}\#T)$. 假定已知 \cref{theorem:ch12-nonhomeomorphic-list} 所列诸紧致曲面是一个完全列表, 那么这个列表中的哪一个曲面同胚于 $P^{2}\#T$?
\end{exercise}

\begin{exercise}
\label{exercise:ch12-sec75-02}
若 $X$ 为 Klein 瓶, 直接计算 $H_{1}(K)$.
\end{exercise}

\begin{exercise}
\label{exercise:ch12-sec75-03}
设 $X$ 为由 $8$-边形区域 $P$ 依标记表 $acadbcb^{-1}d$ 黏合诸边给出的商空间.
\begin{enumerate}[(a)]
\item 验证黏合映射将 $P$ 的所有顶点映射到商空间 $X$ 的一个点.
\item 计算 $H_{1}(X)$.
\item 假设已知 $X$ 同胚于 \cref{theorem:ch12-nonhomeomorphic-list} 中所给出的曲面之一, 这个曲面是哪一个?
\end{enumerate}
\end{exercise}

\begin{exercise}
\label{exercise:ch12-sec75-04}
设 $X$ 是 $8$-边形区域借助于标记表 $abcdad^{-1}cb^{-1}$ 得到的商空间. 设 $\pi\colon P\to X$ 为商映射.
\begin{enumerate}[(a)]
\item 证明 $P$ 的所有顶点在 $\pi$ 下的像不是 $X$ 中的同一个点.
\item 确定空间 $A=\pi(\operatorname{Bd}P)$ 并计算它的基本群.
\item 计算 $\pi_{1}(X,\ x_{0})$ 和 $H_{1}(X)$.
\item 假定已知 $X$ 同胚于 \cref{theorem:ch12-nonhomeomorphic-list} 中给出的曲面之一, 它同胚于哪一个?
\end{enumerate}
\end{exercise}
